\documentclass[12pt,a4paper]{article}
\usepackage{fontspec}
\setmainfont{DejaVu Serif}
\usepackage[russian]{babel}
\usepackage{amsmath,amssymb,mathtools}
\usepackage{geometry}
\usepackage{array,booktabs,longtable}
\usepackage{xcolor}
\usepackage{hyperref}
\usepackage{tikz}
\geometry{left=23mm,right=21mm,top=23mm,bottom=24mm}
\hypersetup{colorlinks=true,linkcolor=blue!55!black,urlcolor=blue!55!black}
\setlength{\parindent}{0pt}
\setlength{\parskip}{0.65em}
\setlength{\emergencystretch}{3em}
\renewcommand{\arraystretch}{1.35}
\begin{document}
\thispagestyle{empty}
\vspace*{0.27\textheight}
\begin{center}
{\LARGE\bfseries Ревью домашней работы\par}
\vspace{2.4cm}
{\large Автор: Лёвин Артём Александрович\par}
\vspace{0.7cm}
{\large 12 октября 2026 года\par}
\end{center}
\vfill
\begin{center}
\begin{tikzpicture}[x=1cm,y=1cm]
\draw[blue!30,line width=0.75pt] (-4.3,0) .. controls (-2.0,0.65) and (-0.7,-0.5) .. (1.15,0.06) .. controls (2.65,0.47) and (3.75,0.16) .. (4.3,0.1);
\end{tikzpicture}
\end{center}
\newpage
\section*{Сводная таблица ошибок}
\footnotesize
\setlength{\tabcolsep}{3pt}
\begin{longtable}{@{}p{0.095\linewidth}p{0.22\linewidth}p{0.27\linewidth}p{0.17\linewidth}p{0.17\linewidth}@{}}
\toprule
\textbf{№} & \textbf{Тема} & \textbf{Суть ошибки} & \textbf{Ваш ответ} & \textbf{Правильный ответ}\\
\midrule
\endfirsthead
\toprule
\textbf{№} & \textbf{Тема} & \textbf{Суть ошибки} & \textbf{Ваш ответ} & \textbf{Правильный ответ}\\
\midrule
\endhead
\hyperref[task:applied3]{3} & Квадратичное неравенство, температурный порог & При делении на отрицательное число знак неравенства не изменён. & $5$ мин & $5$ мин\\
\hyperref[task:applied4]{4} & Подстановка в формулу рейтинга & Вычислен $m$, но не найден рейтинг $R$. & Не указан & $0{,}48$\\
\hyperref[task:applied5]{5} & Кубическая зависимость, радиус сферы & Найден куб радиуса; сам радиус не вычислен. & Не указан & $3{,}9$ м\\
\bottomrule
\end{longtable}
\setlength{\tabcolsep}{6pt}
\normalsize
\emph{Нумерация относится к разделу «Задачи прикладного характера».}
\newpage
\thispagestyle{empty}
\vspace*{\fill}
\begin{center}
{\Huge\bfseries Разбор ошибок}
\end{center}
\vspace*{\fill}
\newpage
\phantomsection\label{task:applied3}
\section*{Задача 3. Нагревательный элемент}
\textbf{Текст задания.} Температура нагревательного элемента (в кельвинах) зависит от времени работы $t$ (в минутах) по формуле
\[
T(t)=T_0+bt+at^2,
\]
где $T_0=680$ К, $b=224$ К/мин, $a=-16$ К/мин$^2$. При температуре свыше $1400$ К прибор может испортиться. Через какое наибольшее время после включения нужно отключить прибор? Ответ дать в минутах.

\textbf{Ваше решение.} Вы записали сравнение температуры с предельно допустимой, перенесли члены в одну часть, а затем разделили выражение на $-16$:
\[
1400>680+224t-16t^2,
\]
\[
-16t^2+224t-720\le 0,
\]
\[
t^2-14t+45\le0.
\]
После этого вычислили дискриминант $D=16$, получили корни $5$ и $9$ и выбрали число $5$.

\textbf{Ваш ответ:} $5$ минут.

\textcolor{red}{\textbf{Ошибка:}} при делении обеих частей неравенства на отрицательное число знак неравенства должен измениться на противоположный. В записи этого не произошло.

\textbf{Где именно возникает ошибка.} Последний верный переход при условии безопасной температуры:
\[
-16t^2+224t-720\le0.
\]
Следующая строка у Вас записана как
\[
t^2-14t+45\le0.
\]
Здесь неверен знак $\le$. Верный результат деления на $-16$:
\[
\boxed{t^2-14t+45\ge0}.
\]

\textbf{Почему этот шаг неверен.} Для неравенств действует правило: если умножить или разделить обе части на отрицательное число, знак сравнения меняется. Например, из $-6\le-2$ после деления на $-2$ получаем $3\ge1$, а не $3\le1$. При делении выражения $-16t^2+224t-720$ на $-16$ коэффициенты меняют знак, и одновременно требуется заменить $\le$ на $\ge$. Ошибка меняет множество допустимых моментов времени. Неравенство $t^2-14t+45\le0$ верно между корнями, то есть как раз на участке превышения температурного порога (кроме граничных точек). Поэтому правильный конечный ответ сам по себе не устраняет ошибку в рассуждении.

\textcolor{blue!60!black}{\textbf{Ключевая идея:}} необходимо найти все моменты, при которых температура \emph{не превышает} $1400$ К, а затем учесть, что прибор нагревается непрерывно с момента включения: нельзя выбрать поздний безопасный момент, если до него прибор уже успел перегреться.

\textbf{Исправление от последнего правильного шага.} Из Вашего неравенства
\[
-16t^2+224t-720\le0
\]
после деления на отрицательное число $-16$ правильно получаем
\[
t^2-14t+45\ge0.
\]
Корни квадратного трёхчлена найдены у Вас верно: $5$ и $9$. Осталось определить, где трёхчлен неотрицателен, и выбрать подходящий промежуток времени.

\textbf{Полное правильное решение.} По условию температура может достигать $1400$ К, но не должна становиться выше этого значения. Следовательно,
\[
680+224t-16t^2\le1400.
\]
Переносим $1400$ влево и приводим подобные слагаемые:
\[
-16t^2+224t-720\le0.
\]
Делим на $-16$, изменив направление знака:
\[
t^2-14t+45\ge0.
\]
Разложим квадратный трёхчлен на множители:
\[
t^2-14t+45=(t-5)(t-9).
\]
Поэтому необходимо решить неравенство
\[
(t-5)(t-9)\ge0.
\]
Произведение двух множителей неотрицательно, когда оба неотрицательны или оба неположительны. Отсюда
\[
t\le5\quad\text{или}\quad t\ge9.
\]
Время работы не может быть отрицательным, поэтому формально допустимые по температуре моменты:
\[
0\le t\le5\quad\text{или}\quad t\ge9.
\]
Но между пятой и девятой минутами температура \emph{больше} допустимой. Так как прибор работает с нулевой минуты без отключения, ждать до девятой минуты нельзя. Максимальная продолжительность безопасной непрерывной работы составляет $5$ минут.

\textbf{Проверка.} На пятой минуте
\[
T(5)=680+224\cdot5-16\cdot25=680+1120-400=1400\ \text{К}.
\]
На шестой минуте
\[
T(6)=680+224\cdot6-16\cdot36=680+1344-576=1448\ \text{К}>1400\ \text{К}.
\]
Следовательно, после пятой минуты порог действительно превышается. При $0\le t\le5$ неравенство безопасности выполняется.

\textcolor{teal!60!black}{\textbf{Итог:}} правильный ответ --- $5$ минут; при этом в решении необходимо исправить знак неравенства и обосновать выбор первого корня.

\textbf{Задача на закрепление.} Температура прибора определяется формулой $T(t)=500+160t-10t^2$, где $t$ --- время в минутах, $T$ --- температура в кельвинах. Если температура выше $1100$ К опасна, через какое наибольшее время после включения необходимо отключить прибор?

\newpage
\phantomsection\label{task:applied4}
\section*{Задача 4. Рейтинг интернет-магазина}
\textbf{Текст задания.} Рейтинг интернет-магазина вычисляется по формулам
\[
R=r_{\text{пок}}-\frac{r_{\text{пок}}-r_{\text{экс}}}{(K+1)^m},
\qquad
m=\frac{0{,}02K}{r_{\text{пок}}+0{,}1},
\]
где $r_{\text{пок}}$ --- средняя оценка покупателей, $r_{\text{экс}}$ --- оценка экспертов, $K$ --- число покупателей, оценивших магазин. Найти рейтинг при $K=15$, $r_{\text{пок}}=0{,}5$, $r_{\text{экс}}=0{,}42$.

\textbf{Ваше решение.} Вы правильно подставили исходные оценки и число покупателей в формулу:
\[
R=0{,}5-\frac{0{,}5-0{,}42}{(15+1)^m}
=0{,}5-\frac{0{,}08}{16^m}.
\]
Затем вычислили показатель степени:
\[
m=\frac{0{,}02\cdot15}{0{,}5+0{,}1}
=\frac{0{,}3}{0{,}6}=0{,}5.
\]
На этом решение закончилось.

\textbf{Ваш ответ:} итоговый рейтинг $R$ не указан (найден только показатель $m=0{,}5$).

\textcolor{red}{\textbf{Ошибка:}} задача решена не до конца. Вспомогательный показатель $m$ найден верно, но требуется вычислить сам рейтинг $R$.

\textbf{Где именно возникает неполнота.} Последний правильный шаг --- $m=0{,}5$. После него нужно вернуться к формуле для $R$ и подставить полученное значение. Этот обязательный шаг отсутствует.

\textbf{Почему это неполное решение.} В задаче присутствуют две величины: вспомогательный показатель степени $m$ и искомый рейтинг $R$. Знание $m$ позволяет вычислить знаменатель, но не является ответом. Здесь показатель $0{,}5$ означает извлечение квадратного корня: $16^{0{,}5}=\sqrt{16}=4$. Только после этого можно вычислить дробь и выполнить вычитание. Остановившись на промежуточном результате, мы не получаем величину, которую просили найти.

\textcolor{blue!60!black}{\textbf{Ключевая идея:}} сначала вычислить вспомогательный показатель $m$, затем использовать его при вычислении степени $(K+1)^m$, после чего найти весь рейтинг $R$.

\textbf{Исправление от последнего правильного шага.} Сохраняем Ваш результат $m=0{,}5$ и подставляем его в уже записанную формулу:
\[
R=0{,}5-\frac{0{,}08}{16^{0{,}5}}.
\]
Поскольку $16^{0{,}5}=4$, получаем
\[
R=0{,}5-\frac{0{,}08}{4}=0{,}5-0{,}02=0{,}48.
\]

\textbf{Полное правильное решение.} Из условия известны $K=15$, $r_{\text{пок}}=0{,}5$ и $r_{\text{экс}}=0{,}42$. Сначала находим показатель:
\[
\begin{aligned}
m&=\frac{0{,}02K}{r_{\text{пок}}+0{,}1}
=\frac{0{,}02\cdot15}{0{,}5+0{,}1}\\
&=\frac{0{,}30}{0{,}60}=0{,}5.
\end{aligned}
\]
Теперь вычисляем числитель дроби:
\[
r_{\text{пок}}-r_{\text{экс}}=0{,}5-0{,}42=0{,}08.
\]
В знаменателе сначала складываем $K+1=16$, затем возводим число $16$ в степень $0{,}5$:
\[
(K+1)^m=16^{0{,}5}=\sqrt{16}=4.
\]
Остаётся подставить эти значения в основную формулу:
\[
\begin{aligned}
R&=r_{\text{пок}}-\frac{r_{\text{пок}}-r_{\text{экс}}}{(K+1)^m}\\
&=0{,}5-\frac{0{,}08}{4}\\
&=0{,}5-0{,}02=0{,}48.
\end{aligned}
\]

\textbf{Проверка.} Числитель поправки положителен и равен $0{,}08$, знаменатель равен $4$, поэтому поправка равна $0{,}02$. Полученное значение $0{,}48$ меньше оценки покупателей $0{,}5$ и больше оценки экспертов $0{,}42$, что согласуется с формулой.

\textcolor{teal!60!black}{\textbf{Итог:}} рейтинг интернет-магазина равен $\boxed{0{,}48}$.

\textbf{Задача на закрепление.} Рейтинг магазина рассчитывается по тем же формулам. Найдите $R$, если $K=24$, $r_{\text{пок}}=0{,}86$ и $r_{\text{экс}}=0{,}61$.

\newpage
\phantomsection\label{task:applied5}
\section*{Задача 5. Максимальный радиус аппарата}
\textbf{Текст задания.} Проектируемый подводный аппарат имеет форму сферы. Выталкивающая сила (в ньютонах) определяется формулой
\[
F_A=\alpha\rho gr^3,
\]
где $\alpha=4{,}2$ --- постоянная, $\rho=1000$ кг/м$^3$ --- плотность воды, $g=10$ Н/кг --- ускорение свободного падения, $r$ --- радиус аппарата в метрах. Найти максимальный радиус, если выталкивающая сила не должна превышать $2\,491\,398$ Н.

\textbf{Ваше решение.} Вы подставили числовые данные:
\[
2\,491\,398=4{,}2\cdot1000\cdot10\cdot r^3,
\]
\[
2\,491\,398=42\,000r^3,
\]
\[
r^3=\frac{2\,491\,398}{42\,000}=59{,}319.
\]
На этом вычисления закончились.

\textbf{Ваш ответ:} радиус в метрах не указан (получено только $r^3=59{,}319$).

\textcolor{red}{\textbf{Ошибка:}} не завершено нахождение искомого радиуса: найден его куб, но не сам радиус. Кроме того, условие «не более» следует записать как неравенство.

\textbf{Где именно возникает неполнота.} Ваше числовое деление выполнено верно. Последний правильный шаг для \emph{граничного} значения силы --- $r^3=59{,}319$. Далее необходимо извлечь кубический корень, получить $r=3{,}9$ и отдельно обосновать, почему именно это значение является максимальным допустимым радиусом.

\textbf{Почему это неполное решение.} Если задано уравнение $r^3=59{,}319$, число $59{,}319$ обозначает \emph{куб радиуса}, а не радиус. Эти величины имеют разные размерности: $r^3$ измеряется в кубических метрах, а в ответе требуются метры. В условии также говорится, что сила не должна превышать заданную границу. Поэтому правильная математическая модель имеет вид $F_A\le2\,491\,398$, а знак равенства возникает лишь при поиске максимального радиуса. Так как для неотрицательного $r$ величина $r^3$ растёт вместе с $r$, большему радиусу соответствует большая выталкивающая сила.

\textcolor{blue!60!black}{\textbf{Ключевая идея:}} для максимального допустимого радиуса выталкивающая сила равна предельному значению; после нахождения $r^3$ нужно извлечь \emph{кубический}, а не квадратный корень.

\textbf{Исправление от последнего правильного шага.} Сохраняем найденное значение
\[
r^3=59{,}319.
\]
Поскольку $3{,}9^3=3{,}9\cdot3{,}9\cdot3{,}9=59{,}319$, получаем
\[
\boxed{r=\sqrt[3]{59{,}319}=3{,}9\ \text{м}}.
\]

\textbf{Полное правильное решение.} По условию
\[
F_A\le2\,491\,398.
\]
Подставим формулу и заданные положительные коэффициенты:
\[
4{,}2\cdot1000\cdot10\cdot r^3\le2\,491\,398.
\]
Умножаем коэффициенты:
\[
42\,000r^3\le2\,491\,398.
\]
Разделим обе части на положительное число $42\,000$, поэтому знак неравенства сохраняется:
\[
r^3\le\frac{2\,491\,398}{42\,000}=59{,}319.
\]
По физическому смыслу радиус неотрицателен: $r\ge0$. При неотрицательных значениях возведение в куб сохраняет порядок чисел. Поэтому
\[
0\le r\le\sqrt[3]{59{,}319}=3{,}9.
\]
Следовательно, максимальный допустимый радиус равен $3{,}9$ м.

\textbf{Проверка.} Подставим найденный радиус в формулу силы:
\[
\begin{aligned}
F_A&=4{,}2\cdot1000\cdot10\cdot(3{,}9)^3\\
&=42\,000\cdot59{,}319=2\,491\,398\ \text{Н}.
\end{aligned}
\]
Сила в точности равна указанному пределу. Любой больший радиус даст больший куб радиуса, а значит, и большую силу, нарушающую условие.

\textcolor{teal!60!black}{\textbf{Итог:}} максимальный радиус аппарата равен $\boxed{3{,}9\ \text{м}}$.

\textbf{Задача на закрепление.} Сферический подводный аппарат испытывает выталкивающую силу $F_A=\alpha\rho gr^3$, где $\alpha=4{,}2$, $\rho=1000$ кг/м$^3$, $g=10$ Н/кг. Найдите максимальный радиус аппарата, если выталкивающая сила должна быть не более $2\,688\,000$ Н. Ответ дайте в метрах.

\vfill
\begin{center}
\small Лёвин Артём Александрович эксклюзивно для Григория Антипова
\end{center}
\end{document}
